\scanpage{88}
\mspage{81}
% archive sheet number: 87
Soit alors $\alpha_0$ le premier ordinal qui soit de cardinal $\deleted{\unsure{\geq|\leq}}>\pi_0$, et considérons \deleted{$\unsure{I|I'}_{\alpha_0}$}
\[
I'=I'_{\alpha_0}=\bigcup_{\substack{\alpha<\alpha_0\\\text{i.e. }\operatorname{card}\alpha\leq\pi_0}}I'_\alpha.
\]
Alors $I'$ contient $I'_0$, il est grand devant $\pi_0$, mais \inserted{quel est} son cardinal. Prouvons par récurrence transfinie que
\[
\operatorname{card}I'_\alpha\leq\pi\qquad
\begin{gathered}
\text{si }\alpha<\alpha_0\\
\text{i.e. }\operatorname{card}\alpha\leq\pi_0.
\end{gathered}
\]
C'est OK pour $\alpha=0$ par hyp. sur $I'_0$. Si c'est vrai pour $I'_\alpha$, est-ce vrai pour $I'_{\alpha+1}$? On \deleted{lui} adjoint à $I'_\alpha$ un ens \unsure{d'éléments|d'élémts} $\unsure{\deleted{e}i_J|i_J}$ qui est de cardinal $\leq\unsure{\operatorname{Card}|\overline{\mathrm C}\mathrm{ard}}\unsure{\mathscr P|\mathcal P|P}_{\leq\pi_0}(I'_\alpha)$ donc de cardinal $\leq\operatorname{card}(I'_\alpha)^{\pi_0}\leq\pi^{\pi_0}$, et si on suppose
\[
\pi^{\pi_0}=\pi,
\]
on \unsure{gagne|gagné}, \deleted{\unsure{Cela sera le cas}} Supposons enfin que $\alpha$ soit un ordinal limite (de cardinal $\leq\pi_0$). Donc on a
\[
I'_\alpha=\bigcup_{\beta<\alpha}I'_\beta
\]
d'où \deleted{$\operatorname{card}$}
